\chapter{Development and Verification}
\label{ch:development}

This chapter describes how the package is checked, what the test suite does
and deliberately does not cover, and the rules a change has to satisfy.  It
is written for someone modifying the package rather than using it.

\section{Running the checks}
\label{sec:dev-verify}

The unified check is one command from the repository root:

\begin{transcript}
$ bash scripts/verify.sh
== syntax gate (compileall) ==
   OK
== test suite ==
........................................................................ [ 88%]
.........                                                                [100%]
81 passed in 0.36s
\end{transcript}

The script runs under \code{set -euo pipefail}, so any failure aborts the
run.  Its first stage is a syntax gate: \code{python -m compileall} over
\file{src/}, \file{scripts/} and \file{tests/}, which fails on any file the
running interpreter cannot parse.  The second stage is the test suite,
\code{python -m pytest -q}, which needs only the core dependencies; no ORCA,
MongoDB, network access or model weights are required, and no external
program is launched.

\section{The test suite}
\label{sec:dev-tests}

\begin{table}[htbp]
\centering
\caption{Collected tests per file.  The counts are collection counts, so a
parametrised test contributes one count per case.}
\label{tab:tests}
\begin{tabularx}{\textwidth}{@{}lrX@{}}
\toprule
File & Tests & Scope \\
\midrule
\file{tests/test_calibration.py} & 24 & the calibration convention and its metrics, pure NumPy \\
\file{tests/test_committee.py} & 21 & the disagreement metrics, pure computation \\
\file{tests/test_factory.py} & 20 & level routing of the calculator factory \\
\file{tests/test_safe_nnp.py} & 16 & the fallback decision rule of the safe calculator \\
\bottomrule
\end{tabularx}
\end{table}

Four design rules keep the suite fast and honest:

\begin{itemize}
  \item \textbf{no external programs.}  Constructing a calculator never runs
        a backend: xTB and ORCA are only invoked by \code{calculate()}, and
        the tests never call it on a real backend.  The committee tests never
        touch a model, because the calculators are built lazily and the
        tests stop before that point.
  \item \textbf{expectations are derived independently.}  The calibration
        tests compare against a hand-written reference quantile (the
        Hyndman--Fan type-7 definition, the NumPy default) and against
        closed-form values written out in the comments, not against numbers
        the production code produced.
  \item \textbf{stubs at the process boundary.}  The safe-calculator tests
        inject a stub committee and a stub fallback, and the production ORCA
        symbol inside the module is monkeypatched, so a regression in the
        decision rule cannot launch a real ORCA process.
  \item \textbf{skip, do not fail, when an optional stack is absent.}  Tests
        that need ASE declare it with \code{pytest.importorskip} at module
        level; a checkout without the optional extras still runs the pure
        NumPy tests.  The root \file{conftest.py} prepends \file{src/} to
        \code{sys.path}, so the suite works without an editable install.
\end{itemize}

\begin{refusalbox}[What the default suite does not cover]
The suite exercises the mathematics and the decision logic; it does not
execute ORCA, xTB, MongoDB, the SCINE engine or real model weights.  Those
paths are covered by the utility scripts of Chapter~\ref{ch:scripts} and by
end-to-end runs (a Puffin job reaching \code{Status.COMPLETE}, a calibrated
threshold on a real reference set).  A green suite is a necessary condition,
not evidence that an engine integration works.
\end{refusalbox}

\section{The Python version floor}
\label{sec:dev-floor}

The package declares \code{requires-python >= 3.11} and classifies for
3.11 and 3.12.  The syntax gate in \code{verify.sh} is what enforces this in
practice: it compiles every source file with the running interpreter, so a
file that uses syntax the interpreter cannot parse fails before the tests
start.  Run the gate with the oldest supported interpreter (3.11) to make it
a floor check; under a newer interpreter it only reflects that interpreter's
own grammar.  This is why the gate is a separate stage rather than a side
effect of importing the package: a syntax error that a newer interpreter
accepts must not reach a user on the floor version.

\section{Bit-for-bit equivalence with upstream}
\label{sec:dev-equivalence}

Files that are adapted from an upstream project -- interface definitions and
protocol implementations rather than new algorithms -- must be shown to be
logically unchanged, so that a modification is a deliberate act rather than
an accident of transcription.  Two complementary checks are used on the
development side:

\begin{itemize}
  \item \textbf{AST comparison with strings masked.}  Parse both files with
        Python's \code{ast} module, replace every string literal (docstrings
        included) by a fixed placeholder, and compare the resulting trees
        with \code{ast.dump}.  A difference in the trees means the logic
        differs; identical trees mean the two files differ only in strings
        and messages.  This is the check that distinguishes ``renamed the
        error message'' from ``changed a comparison''.
  \item \textbf{Normalized code diff.}  Strip comments and docstrings, drop
        blank lines, normalize indentation and whitespace, and diff the two
        texts.  Where the AST comparison gives a yes/no answer, the
        normalized diff gives a readable line-level review of exactly what
        changed.
\end{itemize}

Both checks establish source-level equivalence; neither proves numerical
equivalence.  They are paired with the behavioural evidence -- the test
suite, and where a number is at stake, an independent numerical path -- and
treated as a necessary, not a sufficient, condition.

\section{Contributing}
\label{sec:dev-contributing}

\begin{itemize}
  \item \textbf{The checks must be green.}  \code{bash scripts/verify.sh}
        has to pass before a change is proposed.  A change that cannot pass
        the syntax gate on 3.11 is not ready, whatever it does on a newer
        interpreter.
  \item \textbf{Keep the module boundaries.}  \code{jnuscout.calculators}
        holds ASE calculators and nothing engine-specific;
        \code{jnuscout.scine} holds the bridge and the injection patches;
        \code{jnuscout.uq} holds the pure-NumPy calibration mathematics.
        The top-level import stays cheap: it must not pull in torch, mace or
        scine.
  \item \textbf{Import optional stacks late.}  Heavy or optional
        dependencies are imported inside the function that needs them, and a
        missing stack should raise a clear error at the point of use rather
        than break the package import.  Tests for such code skip when the
        stack is absent.
\end{itemize}
